%-------------------------
% Based off of: https://github.com/sb2nov/resume
% License : MIT
%------------------------

\documentclass[letterpaper,10pt]{article}

\usepackage{latexsym}
\usepackage[empty]{fullpage}
\usepackage{titlesec}
\usepackage{marvosym}
\usepackage[usenames,dvipsnames]{color}
\usepackage{verbatim}
\usepackage{enumitem}
\usepackage[hidelinks]{hyperref}
\usepackage{fancyhdr}
\usepackage[russian, english]{babel}
\usepackage{tabularx}
\input{glyphtounicode}


%----------FONT OPTIONS----------
% sans-serif
% \usepackage[sfdefault]{FiraSans}
% \usepackage[sfdefault]{roboto}
% \usepackage[sfdefault]{noto-sans}
% \usepackage[default]{sourcesanspro}

% serif
% \usepackage{CormorantGaramond}
% \usepackage{charter}


\pagestyle{fancy}
\fancyhf{} % clear all header and footer fields
\fancyfoot{}
\renewcommand{\headrulewidth}{0pt}
\renewcommand{\footrulewidth}{0pt}

% Adjust margins
\addtolength{\oddsidemargin}{-0.5in}
\addtolength{\evensidemargin}{-0.5in}
\addtolength{\textwidth}{1in}
\addtolength{\topmargin}{-.5in}
\addtolength{\textheight}{1.0in}

\urlstyle{same}

\raggedbottom
\raggedright
\setlength{\tabcolsep}{0in}

% Sections formatting
\titleformat{\section}{
	\vspace{-4pt}\scshape\raggedright\large
}{}{0em}{}[\color{black}\titlerule \vspace{-5pt}]

% Ensure that generate pdf is machine readable/ATS parsable
\pdfgentounicode=1

%-------------------------
% Custom commands
\newcommand{\resumeItem}[2]{
	\item\small{
		{\ifx&#1&\else\textbf{#1}: \fi#2 \vspace{-2pt}}
	}
}

\usepackage[T2A]{fontenc}
\usepackage[utf8]{inputenc}
\usepackage{cmap}

\newcommand{\resumeSubheading}[4]{
	\vspace{-1pt}\item
	\begin{tabular*}{0.97\textwidth}{l@{\extracolsep{\fill}}r}
		\textbf{#1} & #2 \\
		\ifx&#3&\else\textit{\small#3} & \textit{\small #4} \\ \fi
	\end{tabular*}\vspace{-5pt}
}


\newcommand{\resumeSubSubheading}[2]{
	\item
	\begin{tabular*}{0.97\textwidth}{l@{\extracolsep{\fill}}r}
		\textit{\small#1} & \textit{\small #2} \\
	\end{tabular*}\vspace{-7pt}
}

\newcommand{\resumeProjectHeading}[2]{
	\item
	\begin{tabular*}{0.97\textwidth}{l@{\extracolsep{\fill}}r}
		\small#1 & #2 \\
	\end{tabular*}\vspace{-7pt}
}

\newcommand{\resumeSubItem}[2]{\resumeItem{#1}{#2}\vspace{-4pt}}

\renewcommand\labelitemii{$\vcenter{\hbox{\tiny$\bullet$}}$}

\newcommand{\resumeSubHeadingListStart}{\begin{itemize}[leftmargin=0.15in, label={}]}
	\newcommand{\resumeSubHeadingListEnd}{\end{itemize}}
\newcommand{\resumeItemListStart}{\begin{itemize}}
	\newcommand{\resumeItemListEnd}{\end{itemize}\vspace{-5pt}}
\newcommand{\ulhref}[2]{\href{#1}{\underline{#2}}}
%-------------------------------------------
%%%%%%  RESUME STARTS HERE  %%%%%%%%%%%%%%%%%%%%%%%%%%%%


\begin{document}

	\begin{center}
		\textbf{\Large \scshape Усатов Павел} \textbf{\Large $|$ Go Backend Developer} \\ \vspace{1pt}
		\small +7 (912) 712-83-52 $|$
		\href{https://t.me/Usatov_Pavel}{\underline{Telegram}} $|$
		\href{mailto:pvusatov@edu.hse.ru}{\underline{pvusatov@edu.hse.ru}} $|$
		\href{https://github.com/UsatovPavel}{\underline{Github}} $|$
		\href{https://usatovpavel.ru}{\underline{\textbf{Сайт}}}
	\end{center}

	%-----------EDUCATION-----------
	\section{Образование}
	\resumeSubHeadingListStart
	\resumeSubheading
	{ВШЭ СПБ}{Санкт-Петербург}
	{Прикладная математика и информатика}{2023-2027}
	\item \textbf{Релевантные курсы:}
	Архитектура компьютера и операционные системы, базы данных, компьютерные сети,
	алгоритмы и структуры данных, архитектура ПО
	\resumeSubHeadingListEnd

	%-----------EXPERIENCE-----------
	\section{Опыт}
	\resumeSubHeadingListStart
	\resumeSubheading
	{Стажировка Backend Т-Банк}{апрель 2026 - июль 2026}
	{Участвовал в миграции документов с long на uuid}{Продуктовая команда документооборота}
	\resumeItemListStart
	\resumeItem{}{Версионировал 10 эндпоинтов под uuid, чтобы не сломать мобильные приложения, установленные у пользователей}
	\resumeItem{}{Написал модуль команды в общем репозитории интеграции: 7 клиентов с документацией взамен legacy трёхлетней давности. Смежные команды завели эпики перехода на новое API}
	\resumeItem{}{Работа с БД: создание алертов, оптимизация запроса, добавление расширений}
	\resumeItem{}{Чинил баги генерации документов: время у нас, восстановил столбец статистики во внешнем сервисе}
	\resumeItemListEnd
	\textbf{Технологии}: Scala, Cats Effect, PostgreSQL, Kafka, Docker, Typescript, Kanban
	\resumeSubheading
	{\ulhref{https://github.com/UsatovPavel/riid}{Утилита RIID для внутреннего облака VK}}{январь 2026 - июнь 2026}
	{Daemon для p2p-загрузки OCI/Docker образов \textit{независимый} от container engine}{Курсовая, научрук - SRE из VK}
	\resumeItemListStart
	\resumeItem{}{Написал \ulhref{https://github.com/UsatovPavel/java-dragonfly-image-puller}{библиотеку}: gRPC-клиент к Dragonfly для загрузки слоёв с соседних нод. \textbf{Взят в использование в VK}}
	\resumeItem{}{Разобрал Go- и Rust-клиенты Dragonfly и обосновал команде переход на Rust вместо согласованной ранее Go: Go не поддерживается с 2026}
	\resumeItem{}{На 20\% быстрее Podman на кластере k8s из 12 нод: 100 образов от 1 МБ до 5 ГБ}
	\resumeItemListEnd
	\textbf{Технологии}: Java, Kubernetes, gRPC, Docker, Podman, OCI/Docker Registry API, p2p, Grafana
	\resumeSubHeadingListEnd

	%-----------PROJECTS-----------
	\section{Проекты}
	\resumeSubHeadingListStart
	\resumeSubheading
	{\ulhref{https://github.com/UsatovPavel/PRAssign}{\textbf{PRAssign}}: Go-сервис управления code review}{ноябрь 2025 - январь 2026}
	{REST API на Gin: команды, пользователи, пул-реквесты, назначение ревьюеров и статистика}{Индивидуальный проект}
	\resumeItemListStart
	\resumeItem{}{Слоистая архитектура cmd/internal: handler $\rightarrow$ service $\rightarrow$ repository на pgx, конфигурация через viper, JWT-аутентификация с проверкой прав автора и админа}
	\resumeItem{}{Асинхронный пайплайн: Go API $\rightarrow$ Kafka (sarama) $\rightarrow$ \ulhref{https://github.com/UsatovPavel/AsyncFactorial}{Scala-consumer} $\rightarrow$ Kafka $\rightarrow$ Go-consumer $\rightarrow$ PostgreSQL, с TTL и фоновой очисткой результатов}
	\resumeItem{}{Nginx балансирует несколько реплик сервиса: повтор запроса по 502/503/504, keepalive, вывод упавшего бэкенда из ротации}
	\resumeItem{}{Пирамида тестов: интеграционные, e2e, smoke и stress; миграции golang-migrate, нагрузка через k6, golangci-lint в пайплайне}
	\resumeItemListEnd
	\textbf{Технологии}: Go, Gin, pgx, sarama, viper, JWT, Kafka, PostgreSQL, Nginx, Docker Compose, k6, golangci-lint

	\resumeSubheading
	{\ulhref{https://github.com/hse-project-Java-2025}{Приложение «TimeTamer»}}{февраль 2025 - август 2025}
	{Серверная часть календаря с AI-ассистентом и геймификацией}{Командный проект}
	\resumeItemListStart
	\resumeItem{}{Общие задания, статистика, достижения, push-уведомления, аутентификация; интеграция ChatGPT и Whisper для голосового ввода}
	\resumeItemListEnd
	\textbf{Технологии}: Java, Spring Boot, PostgreSQL, Docker, JUnit 4/5, Mockito

	\resumeSubHeadingListEnd

	%-----------SKILLS-----------
	\section{Навыки}
	\textbf{Языки}{: Go, Java, Scala} \\
	\textbf{Go}{: Gin, pgx, sarama, viper, golang-migrate, golangci-lint} \\
	\textbf{Технологии и инструменты}{: Docker, PostgreSQL, Kafka, Kubernetes, gRPC, Nginx, k6, Git, Grafana, Prometheus} \\
	\textbf{Навыки}{: REST API, асинхронные пайплайны, тестирование (интеграционное/e2e/нагрузочное), алгоритмы и структуры данных, работа по Kanban}

%-------------------------------------------
\end{document}
